\section{From Operational Cost to Effective Field Operators}
\label{sec:entropic_lagrangian}

The six PM pillars organize the candidate field terms. Their interpretation as a single physical action requires a common carrier state space, update rule, observable map and normalization. We distinguish the proved conditional structures from that remaining identification.

\subsection{Measures, Costs and Evolution}
\label{sec:physics_translation}
\subsubsection{Local densities}
\label{sec:entropic_lagrangian_density}
The local density must be distinguished from its spatial integral and from the accumulated history cost. For a specified spatial measure and time parameter, write
\begin{equation}
 Z_\Phi(t)=\int_{\Sigma_t}\mathcal D(x,t)\,d\mu_t(x),\qquad
 S_\Phi=\int Z_\Phi(t)\,dt.
\end{equation}
This is the spatial-then-temporal organization of the imported action, not an identification of a pointwise density with a system-wide rate. Combining the measures into a covariant spacetime integral requires the physical geometry and its volume element.

A continuum model may replace lattice counting by an integral and decompose a cost into local terms. In four dimensions a dimensionless action has a density of mass dimension four. This dimensional statement does not uniquely specify the terms, their coefficients, or a Lorentzian metric measure. Identifying $d\mu$ with $d^4x\sqrt{-g}$ is part of the proposed continuum gravitational description.

\subsubsection{The statistical and unitary layers}
\label{subsec:observational_filter}
Exponential survival weights define a statistical layer. In the finite-dimensional reconstruction, positive stationary weights and detailed balance make the symmetrized transfer matrix symmetric. Injectivity makes its square positive definite. Its eigenvalues define real logarithmic energies, and multiplication by $e^{-iE_it}$ preserves the amplitude norm (\nolinkurl{Reconstruction}). A deterministic bijection separately preserves that norm by permuting configuration amplitudes.

These are finite-dimensional results under named assumptions. They do not prove that detailed balance follows from zero total flux, that the statistical transfer matrix is injective, or that reconstructed time evolution is a nearest-neighbour update. They do not yet supply a quantum-field reconstruction, the Born rule or measurement dynamics; see \cref{app:quantum_fundamentals}.

\subsection{The Candidate Additive Structure}
\label{subsec:additive_structure}
We organize the proposed action as
\begin{equation}
 \mathcal L_\Phi=\mathcal L_{\rm CAP}+\mathcal L_{\rm MAP}+\mathcal L_{\rm PRO}
 +\mathcal L_{\rm GOV}+\mathcal L_{\rm TOL}+\mathcal L_{\rm MAR}.
\end{equation}
This bookkeeping does not prove statistical independence of six costs or identify six mutually orthogonal physical subspaces. Interactions and shared maintenance must be included without double counting. Costs incurred by coordination and savings relative to an uncoordinated baseline are distinct entries: a negative contribution to a net comparison is not a negative primitive expenditure. A Lorentzian Lagrangian need not be pointwise positive, so this accounting distinction does not make every displayed field term a positive dissipation density.

\subsection{Kinematic Constraints}
\label{subsec:kinetic_constraint}
A nearest-neighbour update with spacing $\ell$ and step $\tau$ has a specified causal stencil. Choosing units $\ell/\tau=1$ sets its link-speed scale; it does not by itself derive an isotropic relativistic dispersion, Lorentz transformations or equality of all field speeds. The walk calculations exhibit why this matters: the one-coin update has twins, while a particular twin-gapping split-step family reduces its light speed. The two-sector Hamiltonian is a different candidate. A shared low-energy metric requires matching these responses.

\subsection{Dimensions and the Leading Operator List}
\label{subsec:mapping_rules}
\label{sec:equipartition}
For the stipulated first-order fermion and second-order boson kinetic terms, dimensional consistency gives
\begin{equation}
 [\psi]=3/2,\qquad [\phi]=[A]=1.
\end{equation}
Lean proves these relations and constraints on monomials of dimension at most four (\nolinkurl{Operators}). In that list a fermion bilinear admits at most one additional unit of dimension, Yukawa and quartic scalar terms are allowed, and four-fermion and sixth-order scalar terms lie outside the cutoff.

The cutoff is a hypothesis defining the leading operator list. Finite state count alone does not prove it as an exact exclusion: \nolinkurl{Potential.finite_states_do_not_force_quartic} distinguishes a sixth-order interaction on four amplitudes from every proposed quadratic-plus-quartic potential. With coefficients $c\ell^k$, higher-dimensional terms vanish as $\ell\to0$ (\nolinkurl{Operators.higher_dimension_suppressed}); they need not vanish at finite spacing.

\subsubsection{Mass terms and potentials}
\label{subsec:restmass}
The leading list permits a dimension-one fermion mass coefficient and scalar quadratic/quartic terms. Representation and gauge-covariance constraints must decide whether a bare mass or a scalar-mediated coupling is allowed. Orthogonality of two geometric sectors alone is not that representation-theoretic selection rule.

\subsection{Target Continuum Form}
\label{subsec:the_target}
The program seeks a response of the form
\begin{equation}
\begin{aligned}
 \mathcal L_\Phi={}&\bar\psi(i\gamma^\mu D_\mu-m)\psi
 -\tfrac14 F^a_{\mu\nu}F^{a\mu\nu}
 -|D_\mu\phi|^2-V(\phi)\\
 &+\tfrac{M_P^2}{2}(R-2\Lambda).
\end{aligned}
\label{eq:entropic_lagrangian}
\end{equation}
This is a target effective description, with consistent metric and field conventions required. The following sections prove parts of its operator structure for specified models. They do not yet derive its full representation content, gravitational dynamics or coefficients from the substrate.
